% One-time imported draft. Editable native source; no Markdown regeneration.
\section{A board evolves by local
rules}\label{008-life-generations-p00-a-board-evolves-by-local-rules}

A cellular automaton is an array of cells updated by a common local
rule. In Life, each cell is dead or alive. The next state depends on
that cell and the eight positions surrounding it. The same rule is
applied throughout a generation, producing a new board.

Conway\textquotesingle s Life was presented in Martin
Gardner\textquotesingle s October 1970
\href{https://www.scientificamerican.com/article/mathematical-games-1970-10/}{Scientific
American column}. The rule used here is often written B3/S23: birth with
three neighbors, survival with two or three. The chapter uses that rule
on a deliberately finite board with a permanently dead exterior.
Boundary behavior is part of this project, not an implicit property of
Life.

The default pattern is a horizontal line of three cells. After one step
it is vertical; after another it is horizontal again. This small
oscillator makes a repeated final picture possible even when an
implementation has skipped or corrupted the intermediate work. The
program will therefore expose every complete generation, not just its
last image.
